% Section: frames

\documentclass[../../../include/open-logic-section]{subfiles}

\begin{document}

\olfileid[fr]{nml}{frd}{fra}

\olsection{Cadres}

\begin{defn}
  Un \emph{cadre} est une paire $\mModel{F} = \tuple{W,R}$,
  où $W$ est un ensemble non vide de mondes et $R$ une relation
  binaire sur~$W$. Un modèle $\mModel{M}$ est \emph{fondé sur}
  le cadre $\mModel{F} = \tuple{W,R}$ si et seulement si
  $\mModel{M} = \tuple{W, R, V}$ pour une certaine attribution~$V$.
\end{defn}

\begin{defn}
  Si $\mModel{F}$ est un cadre, nous disons que $!A$ est
  \emph{valide dans~$\mModel{F}$}, ce que nous notons
  $\mModel{F} \Entails !A$, si $\mSat{M}{!A}$ pour tout
  modèle~$\mModel{M}$ fondé sur~$\mModel{F}$.

  Si $\mClass{F}$ est une classe de cadres, nous disons que
  $!A$ est \emph{valide dans~$\mClass{F}$}, ce que nous notons
  $\mClass{F} \Entails !A$, si et seulement si
  $\mModel{F} \Entails !A$ pour tout
  cadre~$\mModel{F} \in \mClass{F}$.
\end{defn}

La vérité d'un schéma dans un modèle porte sur la valuation
de ce modèle. Sa validité dans un cadre porte sur toutes
les valuations sur ce cadre. Cette quantification permet d'établir
les correspondances entre schémas et propriétés de la relation
d'accessibilité~$R$. Ainsi, le schéma \Ax{T} est valide dans un
cadre si et seulement si le cadre est réflexif. Il est donc vrai
dans tous les modèles réflexifs, et il peut aussi être vrai dans
un modèle non réflexif. Prenons deux mondes dont chacun a pour
seul successeur l'autre, et donnons à chaque variable
propositionnelle la même valeur dans les deux mondes. Par
induction sur les formules, toute formule a la même valeur dans
les deux mondes, puisque leurs successeurs s'échangent.
Ainsi, toute instance de \Ax{T} est vraie dans ce modèle.

\begin{rem}
La validité dans une classe de cadres est un cas particulier
de la validité dans une classe de modèles :
$\mClass{F} \Entails !A$ si et seulement si
$\mClass{C} \Entails !A$, où $\mClass{C}$ est la classe
de tous les modèles fondés sur un cadre de~$\mClass{F}$.

Si une !!{formula} ou un schéma est valide, c'est-à-dire
valide relativement à la classe de \emph{tous} les modèles,
il est aussi valide relativement à toute classe~$\mClass{F}$
de cadres.
\end{rem}

\end{document}
